%! Author = nathanwhite
%! Date = 8/29/26

% Preamble
\documentclass[11pt]{article}

% Packages
\usepackage[margin=1in]{geometry}
\usepackage{amsmath,amssymb,bm}
\usepackage[hidelinks]{hyperref}
\usepackage[numbers,sort&compress]{natbib}

% Document
\begin{document}

\title{A Multiscale CMC-Fracture Problem:\\Motivation, Mathematical Statement, and Methodological Basis}
\author{}
\date{}
\maketitle

\begin{abstract}
Ceramic-matrix composites (CMCs) offer low density, high-temperature capability,
and damage tolerance, but their matrix-cracking response depends on geometry,
material architecture, interfaces, loading, and environment. This note defines
a bounded computational problem: evaluate a fracture-driving quantity and its
associated fields with a finite-element (FE) model, learn a fast surrogate for
the same parameterized map, and accept its use only when it agrees with the
reference calculation under predeclared criteria.
\end{abstract}

\section{Motivation}
Continuous-fiber CMCs are intended to avoid the abrupt tensile failure of
monolithic ceramics.  A suitably compliant interphase can deflect matrix cracks,
permit interface debonding, and leave fibers bridging the crack faces.  The
material can then accumulate distributed damage rather than concentrating all
energy at one unstable crack.  The compensating difficulty is that the response
depends on constituent properties, fiber architecture, interface condition,
defects, temperature history, and the local multiaxial stress field.  NASA
reports on SiC/SiC composites identify matrix cracking, oxidation, creep, and
fiber degradation as coupled life-relevant phenomena
\citep{nasacmcbehavior,nasasic1200}.

For an analyst, a CAD geometry is consequently not a material model.  It must be
paired with a declared architecture, constitutive and fracture laws, boundary
conditions, loading history, and mesh.  Its results must retain units,
discretization, solver settings, and material parameters.  Otherwise two outputs
called ``fracture'' can answer different physical questions.  This distinction is
unromantic, but a simulation pipeline does not become credible simply because its
containers are arranged in a directed acyclic graph.

The problem here is deliberately narrower than complete service-life prediction.
For a family of CMC structures, estimate displacement and stress fields near a
crack or notch and a fracture-driving quantity represented by the $J$-integral.
In its original elastic setting, $J$ is a path-independent measure of energy
release rate for crack advance \citep{rice1968}.  In a heterogeneous,
bridged, nonlinear, or dissipative CMC, path independence and equivalence to a
unique material toughness require additional assumptions.  Thus $J$ is an
informative target and diagnostic, not a scalar that makes the microphysics go
away.

High-fidelity FE analysis gives a defensible reference route: it resolves the
specified geometry and boundary conditions, exposes constitutive assumptions,
and permits mesh and contour-convergence studies.  Its cost becomes onerous in a
study spanning geometries, thermal states, material parameters, and loads.  A
neural operator is attractive because it learns a map between input functions and
solution fields, rather than just a finite-dimensional regression at one grid
resolution \citep{kovachki2023}.  The Fourier neural operator (FNO) uses
Fourier-space kernel parameterization and has demonstrated rapid PDE-solution
surrogates on benchmark families \citep{li2021fno}.

Speed is not evidence.  A model trained on FE cases can interpolate within its
corpus yet fail on a changed notch topology, fiber orientation, thermal state, or
crack configuration.  Fracture quantities are especially unforgiving: a modest
local field error can change a contour integral or a threshold decision.  The FE
track must therefore remain a comparison authority for selected cases.  The
system should record disagreement, quality checks, and an out-of-distribution
indicator, then refuse a surrogate screening result when the stated limits are
not met.  That refusal is a useful result, not a software exception wearing a
false moustache.

This has consequences for how the multiscale problem is framed.  The word
``multiscale'' should not suggest that a component calculation has automatically
resolved the fiber, interface, pore, matrix grain, and crack process zone at
their physical scales.  It means that information from more than one scale is
represented deliberately.  A component-scale FE model may use homogenized
anisotropic properties, an interface law informed by smaller-scale experiments,
and a crack-driving functional evaluated at the structural scale.  Each
abstraction has a regime of validity.  In particular, the parameters in an
effective law should be regarded as calibrated model inputs, not as immutable
properties that remain correct after changes in temperature, environmental
exposure, architecture, or manufacturing route.

The scale separation also explains why the choice of output matters.  A scalar
such as $J$ is valuable because it supports comparison across a parameter sweep,
but it is not sufficient to diagnose every failure mechanism.  A result viewer
should preserve the displacement, stress, and damage fields that explain why a
particular value was obtained.  A high $J$ near one region may be caused by a
geometric stress concentrator, a boundary-condition artifact, an inadequate mesh,
or a genuine material-interface effect.  Conflating those possibilities into a
single traffic-light label would be efficient in the same sense that a compass is
an efficient substitute for a map.

The proposed FE--surrogate pairing is therefore best understood as an
experiment-design mechanism.  The expensive model supplies selected reference
cases across the variables that define the intended domain of use.  The fast
model explores between those cases and highlights candidates that merit renewed
resolution.  This permits adaptive allocation of computational effort, but it
also requires discipline in data splitting and reporting.  Randomly withholding
individual voxels or repeated perturbations of one geometry is not evidence that
a surrogate can generalize to a new component family.  The held-out cases should
match the actual claim: unseen loads for load interpolation, unseen geometries
for geometric generalization, or unseen material states for material
generalization.  These are different questions and usually have different error
profiles.

Finally, the workflow has a useful engineering boundary.  Its near-term purpose
is comparative screening and model development, not automated release of a
fracture-critical design.  To cross that boundary would require model-form
verification, calibration, experimental validation in the relevant regime,
uncertainty treatment, and configuration control of both numerical and learned
models.  Those activities are expensive because they replace attractive diagrams
with accountable evidence.  They are also the activities that make a prediction
worth believing.  The formulation below is intentionally arranged so that this
later evidence can attach to explicit inputs, outputs, assumptions, and
acceptance decisions.

\section{Problem to be solved}

Let $\Omega(\theta)\subset\mathbb{R}^3$ be a component domain obtained from
geometry parameters $\theta$.  Let $\Gamma_u$, $\Gamma_t$, and $\Gamma_c$
denote displacement, traction, and crack/damage boundaries.  A simulation case
is
\begin{equation}
 x=\bigl(\Omega(\theta),\mathcal{A},\mathcal{M},b,\ell,\tau\bigr),
\end{equation}

where $\mathcal{A}$ is the modeled material architecture, $\mathcal{M}$ its
constitutive/fracture parameters, $b$ boundary data, $\ell$ a load history, and
$\tau$ the analysis increment.

Its required output is

\begin{equation}
 y=\bigl(\bm{u},\bm{\sigma},\bm{\varepsilon},d,J;q\bigr),
\end{equation}

where the first four terms are displacement, stress, strain, and a declared
damage/crack state; $J$ is a fracture quantity; and $q$ records provenance,
units, mesh, model/solver versions, tolerances, and convergence status.

For a quasistatic reference model, seek $\bm{u}\in\mathcal{V}$ such that
\begin{equation}
 \int_{\Omega}\bm{\varepsilon}(\bm{v}):
 \bm{\sigma}(\bm{u},d,\mathcal{M})\,d\Omega
 =\int_{\Omega}\bm{v}\cdot\bm{f}\,d\Omega+
 \int_{\Gamma_t}\bm{v}\cdot\bar{\bm{t}}\,d\Gamma
 \quad\forall\bm{v}\in\mathcal{V}_0,
 \label{eq:equilibrium}
\end{equation}
with $\bm{u}=\bar{\bm{u}}$ on $\Gamma_u$.  This equilibrium statement is not
yet a complete CMC model.  An implementation must state its kinematics,
anisotropy, thermal dependence, interface law, and damage evolution.  For
example, an energy-based damage model enforces $\dot d\geq0$; a cohesive model
instead introduces a traction--separation law on an internal surface.  They are
alternative hypotheses, not interchangeable labels.

For a contour $\Gamma$ surrounding a crack tip with local propagation direction
$x_1$, write
\begin{equation}
 J=\int_{\Gamma}\left(Wn_1-\sigma_{ij}n_j
 \frac{\partial u_i}{\partial x_1}\right)ds,
 \qquad W=\int_0^{\bm{\varepsilon}}\bm{\sigma}:d\bm{\varepsilon}.
 \label{eq:j}
\end{equation}
In the CMC application, every reported value must identify contour/domain,
crack-front direction, treatment of inhomogeneity and dissipation, and mesh and
contour convergence.  This is what makes $J$ auditable rather than merely
numerically expensive.

Denote the reference operator by $\mathcal{S}:x\mapsto y_{\mathrm{FE}}$.  Given
a corpus $\mathcal{D}_{\mathrm{train}}$ of verified reference cases, train an
FNO surrogate $\widehat{\mathcal{S}}_\phi$ to approximate selected fields and
functionals.  A schematic FNO layer is
\begin{equation}
 \bm{v}_{k+1}(\bm{r})=\sigma\!\left(W_k\bm{v}_k(\bm{r})+
 \mathcal{F}^{-1}\!\left(R_k\mathcal{F}[\bm{v}_k]\right)(\bm{r})\right),
\end{equation}
with local mixing $W_k$, retained Fourier modes $R_k$, and transform
$\mathcal{F}$.  Inputs must have a declared common raster or geometry-aware
encoding; FFT-based processing does not make an arbitrary unstructured FE mesh
into a regular voxel field.

The training loss should weight the outputs used in the actual decision:
\begin{equation}
 \min_\phi\;\mathbb{E}_{(x,y)\sim\mathcal{D}_{\mathrm{train}}}
 \left[\alpha\frac{\|\widehat{\bm{u}}-\bm{u}\|_2^2}{\|\bm{u}\|_2^2+\epsilon}
 +\beta\frac{\|\widehat{\bm{\sigma}}-\bm{\sigma}\|_2^2}{\|\bm{\sigma}\|_2^2+\epsilon}
 +\gamma\frac{|\widehat{J}-J|}{|J|+\epsilon}\right].
 \label{eq:loss}
\end{equation}
Held-out evaluation must separate the physical generalization claim---new loads,
geometries, or material states---rather than neighboring records from the same
run.  A new case is accepted for screening only when field and $J$ errors,
numerical-quality checks, and an out-of-distribution criterion satisfy
predeclared limits.  Otherwise it is routed to FE and labeled \emph{indeterminate}.

\section{Brief literature review and methodological position}

The mechanics foundation begins with energy-based fracture measures.  Rice
introduced the path-independent integral behind Equation~\eqref{eq:j}
\citep{rice1968}.  For CMCs, Marshall and Evans documented multiple matrix
cracking, fiber fracture, and pull-out, while Marshall, Cox, and Evans modeled
the role of fiber-bridging closure tractions in matrix cracking
\citep{marshallEvans1985,marshallCoxEvans1985}.  Budiansky, Hutchinson, and
Evans distinguish interface regimes relevant to matrix fracture
\citep{budiansky1986}; Moran and Shih provide the domain-integral formulation
used naturally in finite-element evaluation \citep{moranShih1987}.  Together,
this literature argues for explicit crack/interface assumptions rather than a
single, context-free toughness number.

FE is the reference mechanism here because it makes governing assumptions and
numerical sensitivity inspectable.  Reference does not mean physical truth:
calibration, experimental validation, and uncertainty evidence are still needed
before any result supports an allowable or life claim.  The surrogate literature
offers a complementary route.  Neural operators formalize maps between function
spaces \citep{kovachki2023}; FNOs are a Fourier-parameterized practical member
of that family \citep{li2021fno}.  Geo-FNO extends the approach toward irregular
geometries \citep{li2022geofno}.  Neither result establishes accuracy for CMC
crack fields or $J$; that must be measured against the declared FE corpus.

The methodological position is therefore conservative.  The surrogate is an
acceleration layer with a bounded empirical domain of use.  It can prioritize
designs and make field-level exploration cheaper, but it may not silently
replace the reference calculation.  Its comparison stage is an adjudication
record, not validation unless independent physical evidence has been supplied.

\section{Scope and next evidence}

The immediate scientific deliverable is one reproducible benchmark corpus for a
specified CMC architecture and fracture configuration.  Every case should carry
the input contract, reference fields, $J$ evaluation details, mesh/convergence
evidence, and stable artifact identity.

Only then can accuracy, speedup, and
refusal behavior be measured.

Temperature dependence, interface degradation,
transient loads, and richer microstructure are appropriate later extensions, but
each expands the physical claim and requires fresh verification and validation.

\bibliographystyle{plainnat}
\bibliography{theory}

\end{document}
